% ---------------------------------------------------------------
% preambulo.tex — pacotes, estilos e ambientes
% ---------------------------------------------------------------

\usepackage[T1]{fontenc}
\usepackage[utf8]{inputenc}
\usepackage[portuguese]{babel}
\usepackage{lmodern}
% A bibliografia chama-se «Referências», como na versão Word do relatório.
\addto\captionsportuguese{\renewcommand{\bibname}{Referências}}
\usepackage{microtype}

\usepackage[a4paper,margin=2.6cm,headheight=15pt]{geometry}
\usepackage{graphicx}
\usepackage{booktabs}
\usepackage{longtable}
\usepackage{array}
\usepackage{tabularx}
\usepackage{multirow}
\usepackage{amsmath}
\usepackage{amssymb}
\usepackage{enumitem}
\usepackage[table]{xcolor}
\usepackage{fancyhdr}
\usepackage{titlesec}
\usepackage{caption}
\usepackage{float}
\usepackage{listings}
% Bibliografia no índice, com a página e a ligação corretas (no início da
% bibliografia), sem acrescentar o próprio índice nem as listas.
\usepackage[nottoc,notlof,notlot]{tocbibind}
\usepackage{csquotes}
% Aspas: \enquote{texto} produz “texto” (aspas duplas tipográficas).
\setquotestyle{english}
% O babel-portuguese torna " um atalho que «come» o carácter seguinte;
% desativa-se, para que aspas escritas diretamente não estraguem o texto.
\AtBeginDocument{\shorthandoff{"}}
\usepackage{icomma}        % vírgula decimal em modo matemático: $1,5$
\usepackage[most]{tcolorbox}
\usetikzlibrary{positioning}  % diagramas TikZ (o tcolorbox já carrega o tikz)

% --- Paleta ----------------------------------------------------
\definecolor{corEscura}{HTML}{1F3A5F}
\definecolor{corDestaque}{HTML}{AD451D}  % terracota do logo do MEIA
\definecolor{corSucesso}{HTML}{0B6E4F}   % verde: apenas marcadores PASSOU
\definecolor{corCinza}{HTML}{5A6672}
\definecolor{corFundo}{HTML}{F4F6F8}
\definecolor{corCritico}{HTML}{9B1C1C}
\definecolor{corAviso}{HTML}{B45309}

% obeyspaces/spaces: \cod{swipl -q -g main} mantém os espaços e pode quebrar neles.
\usepackage[obeyspaces,spaces]{url}
\DeclareUrlCommand\cod{\urlstyle{tt}}

\usepackage[pdfusetitle]{hyperref}
\hypersetup{
  colorlinks=true,
  linkcolor=corEscura,
  citecolor=corEscura,
  urlcolor=corDestaque,
  linktoc=all
}

% O cleveref não traz nomes em português; definem-se aqui.
\usepackage[capitalise,noabbrev,nameinlink]{cleveref}
\crefname{chapter}{Capítulo}{Capítulos}
\crefname{section}{Secção}{Secções}
\crefname{subsection}{Secção}{Secções}
\crefname{appendix}{Anexo}{Anexos}
\crefname{figure}{Figura}{Figuras}
\crefname{table}{Tabela}{Tabelas}
\crefname{equation}{Equação}{Equações}
\crefname{enumi}{Item}{Itens}
\newcommand{\crefrangeconjunction}{ a }
\newcommand{\crefpairconjunction}{ e }
\newcommand{\crefmiddleconjunction}{, }
\newcommand{\creflastconjunction}{ e }
% O cleveref usa o tipo «listing» para o contador lstlisting.
\crefname{listing}{Listagem}{Listagens}

% --- Cabeçalhos e rodapés --------------------------------------
\pagestyle{fancy}
\fancyhf{}
\fancyhead[L]{\small\textcolor{corCinza}{\ucsigla{} --- \cursosigla}}
\fancyhead[R]{\small\textcolor{corCinza}{\tituloprincipal}}
\fancyfoot[C]{\thepage}
\renewcommand{\headrulewidth}{0.4pt}

\fancypagestyle{plain}{%
  \fancyhf{}%
  \fancyfoot[C]{\thepage}%
  \renewcommand{\headrulewidth}{0pt}%
}

% --- Títulos ---------------------------------------------------
\titleformat{\chapter}[hang]
  {\normalfont\huge\bfseries\color{corEscura}}
  {\thechapter\hspace{12pt}\textcolor{corCinza}{|}\hspace{12pt}}{0pt}{\huge}
\titlespacing*{\chapter}{0pt}{-10pt}{24pt}

\titleformat{\section}
  {\normalfont\large\bfseries\color{corEscura}}{\thesection}{0.8em}{}
\titleformat{\subsection}
  {\normalfont\normalsize\bfseries\color{corEscura}}{\thesubsection}{0.8em}{}

% Uma secção não começa no fundo da página: se restar pouco espaço,
% passa para a página seguinte (no máximo ~1/5 da página fica em branco).
\usepackage{needspace}
\usepackage{placeins}  % \FloatBarrier: impede uma figura de passar para a secção seguinte
\usepackage{etoolbox}
\pretocmd{\section}{\Needspace{10\baselineskip}}{}{}
\pretocmd{\subsection}{\Needspace{6\baselineskip}}{}{}

% --- Legendas --------------------------------------------------
\captionsetup{font=small,labelfont={bf,color=corEscura},skip=6pt}

% --- Listagens de código ---------------------------------------
% Drools (DRL): sintaxe próxima do Java, com as palavras-chave das regras.
% Usada apenas na comparação com a regra Prolog equivalente (language=Drools).
\lstdefinelanguage{Drools}{
  morekeywords={package,import,global,function,declare,query,rule,when,
                then,end,salience,agenda-group,activation-group,no-loop,
                lock-on-active,dialect,and,or,not,exists,forall,from,
                collect,accumulate,eval,insert,insertLogical,update,modify,
                retract,delete,new,true,false,null,this,matches,memberOf,
                contains,extends,if,else,return},
  otherkeywords={agenda-group,activation-group,no-loop,lock-on-active},
  sensitive=true,
  comment=[l]{//},
  morecomment=[n]{/*}{*/},
  string=[b]",
  morestring=[b]'
}

\lstdefinestyle{codigo}{
  basicstyle=\ttfamily\footnotesize,
  keywordstyle=\color{corEscura}\bfseries,
  commentstyle=\color{corCinza}\itshape,
  stringstyle=\color{corDestaque},
  numbers=left,
  numberstyle=\tiny\color{corCinza},
  numbersep=8pt,
  backgroundcolor=\color{corFundo},
  frame=leftline,
  framerule=1.2pt,
  rulecolor=\color{corEscura},
  framesep=8pt,
  xleftmargin=16pt,
  breaklines=true,
  breakatwhitespace=true,
  showstringspaces=false,
  columns=fullflexible,
  captionpos=b,
  language=Drools,
  literate=%
    {á}{{\'a}}1 {à}{{\`a}}1 {â}{{\^a}}1 {ã}{{\~a}}1
    {é}{{\'e}}1 {ê}{{\^e}}1 {í}{{\'i}}1
    {ó}{{\'o}}1 {ô}{{\^o}}1 {õ}{{\~o}}1
    {ú}{{\'u}}1 {ç}{{\c c}}1
    {Á}{{\'A}}1 {É}{{\'E}}1 {Í}{{\'I}}1 {Ó}{{\'O}}1 {Ú}{{\'U}}1 {Ç}{{\c C}}1
}
% Prolog (linguagem por omissão neste relatório) e JSON (pedidos ao server.pl).
% No Prolog, o apóstrofo delimita átomos ('Nome da regra'); as aspas, strings.
\lstdefinelanguage{JSON}{
  string=[b]",
  sensitive=true,
  morekeywords={true,false,null}
}
\lstset{style=codigo,language=Prolog}

\renewcommand{\lstlistingname}{Listagem}
\renewcommand{\lstlistlistingname}{Lista de Listagens}

% --- Macros auxiliares -----------------------------------------
\newcommand{\facto}[1]{\texttt{\small #1}}
\newcommand{\regra}[1]{\texttt{\small #1}}
% Identificadores de código no texto (predicados, ficheiros, classes), com
% \cod{...}: definido mais abaixo com o pacote url, que dispensa escapar os _
% de nomes como wm_support/3 e permite quebrar a linha depois de _ / . e nos
% espaços. Só ASCII dentro de \cod (para texto acentuado, usar \texttt com \_).
\newcommand{\pts}[1]{\textbf{#1}}

% Informação em falta (a obter do perito ou dos autores). Fica realçada
% para não passar despercebida para a versão final.
\newcommand{\preencher}[1]{\colorbox{corAviso!20}{\textcolor{corAviso}{\small\textbf{[#1]}}}}

% Regra em linguagem natural (SE ... ENTÃO ...), com o identificador no título.
% Uso: \begin{regranl}{R_CHUVA_01} SE ... \end{regranl}
\newtcolorbox{regranl}[1]{%
  enhanced, breakable,
  colback=corFundo, colframe=corEscura,
  boxrule=0pt, leftrule=1.2pt, arc=0pt, outer arc=0pt,
  left=8pt, right=8pt, top=4pt, bottom=4pt,
  fonttitle=\ttfamily\small\bfseries, coltitle=corEscura,
  attach title to upper, after title={\par\smallskip},
  title={REGRA \detokenize{#1}},
  fontupper=\ttfamily\small
}

\newcolumntype{L}[1]{>{\raggedright\arraybackslash}p{#1}}
\newcolumntype{C}[1]{>{\centering\arraybackslash}p{#1}}

% --- Grelha das tabelas -----------------------------------------
% V: linha vertical fina e cinzenta entre colunas, p. ex.
%    \begin{tabularx}{\textwidth}{@{}L{3cm}VX@{}}
% \linhasep: linha horizontal clara entre linhas de dados, a usar
%    depois de \\ (exceto na última linha, antes do \bottomrule).
\newcolumntype{V}{!{\color{corCinza}\vrule width 0.4pt}}
\newcommand{\linhasep}{%
  \arrayrulecolor{corCinza!45}\specialrule{0.3pt}{0pt}{0pt}%
  \arrayrulecolor{black}}
% Sem o espaço que o booktabs deixa junto às regras, para que as linhas
% verticais fiquem contínuas; o \extrarowheight compensa esse espaço.
% Aplica-se só às tabelas de conteúdo (a capa usa tabular).
\newcommand{\grelhatabelas}{%
  \setlength{\aboverulesep}{0pt}%
  \setlength{\belowrulesep}{0pt}%
  \setlength{\extrarowheight}{2pt}}
\AtBeginEnvironment{tabularx}{\grelhatabelas}
\AtBeginEnvironment{longtable}{\grelhatabelas}

% Margem extra que o TeX pode usar como ultimo recurso antes de deixar
% uma linha a transbordar. Evita \emph{overfull boxes} pontuais em
% paragrafos com termos longos em teletipo (nomes de predicados).
\setlength{\emergencystretch}{3em}

\setlist{itemsep=2pt,parsep=0pt,topsep=4pt}
\renewcommand{\arraystretch}{1.25}
